% Bagian: normal-modal-logic
% Bab: tableaux
% Subbagian: more-rules

\documentclass[../../../include/open-logic-section]{subfiles}

\begin{document}

\olfileid[id]{nml}{tab}{mru}

\olsection{Aturan bagi Relasi Aksesibilitas Lain}

Untuk menangani logika yang ditentukan oleh relasi aksesibilitas khusus,
kita membahas aturan tambahan dalam \olref{tab:more-rules}.

\begin{table}
  \begin{center}
    \def\arraystretch{3}\def\fCenter{}
    \iftag{notprvBox,notprvDiamond}
          {\begin{tabular}{|c|}}
          {\begin{tabular}{|c|c|}}
    \hline
    \iftag{prvBox}{
      \AxiomC{\sFmla{\True}{\Box !A}[\sigma]}
      \RightLabel{T$\Box$}
      \UnaryInfC{\sFmla{\True}{!A}[\sigma]}
      \DisplayProof}{}
    \iftag{notprvBox,notprvDiamond}{}{&}
    \iftag{prvDiamond}{
      \AxiomC{\sFmla{\False}{\Diamond !A}[\sigma]}
      \RightLabel{T$\Diamond$}
      \UnaryInfC{\sFmla{\False}{!A}[\sigma]}
      \DisplayProof}{}
    \\[1ex]
    \hline
    \iftag{prvBox}{
      \AxiomC{\sFmla{\True}{\Box !A}[\sigma]}
      \RightLabel{D$\Box$}
      \UnaryInfC{\sFmla{\True}{\Diamond!A}[\sigma]}
      \DisplayProof}{}
    \iftag{notprvBox,notprvDiamond}{}{&}
    \iftag{prvDiamond}{
      \AxiomC{\sFmla{\False}{\Diamond !A}[\sigma]}
      \RightLabel{D$\Diamond$}
      \UnaryInfC{\sFmla{\False}{\Box!A}[\sigma]}
      \DisplayProof}{}
    \\[1ex]
    \hline
    \iftag{prvBox}{
      \AxiomC{\sFmla{\True}{\Box !A}[\sigma.n]}
      \RightLabel{B$\Box$}
      \UnaryInfC{\sFmla{\True}{!A}[\sigma]}
      \DisplayProof}{}
    \iftag{notprvBox,notprvDiamond}{}{&}
    \iftag{prvDiamond}{
      \AxiomC{\sFmla{\False}{\Diamond !A}[\sigma.n]}
      \RightLabel{B$\Diamond$}
      \UnaryInfC{\sFmla{\False}{!A}[\sigma]}
      \DisplayProof}{}
    \\[1ex]
    \hline
    \iftag{prvBox}{
      \AxiomC{\sFmla{\True}{\Box !A}[\sigma]}
      \RightLabel{4$\Box$}
      \UnaryInfC{\sFmla{\True}{\Box!A}[\sigma.n]}
      \DisplayProof}{}
    \iftag{notprvBox,notprvDiamond}{}{&}
    \iftag{prvDiamond}{
      \AxiomC{\sFmla{\False}{\Diamond !A}[\sigma]}
      \RightLabel{4$\Diamond$}
      \UnaryInfC{\sFmla{\False}{\Diamond!A}[\sigma.n]}
      \DisplayProof}{}
    \\
    \iftag{notprvBox,notprvDiamond}
         {$\sigma.n$ sudah digunakan}
         {$\sigma.n$ sudah digunakan & $\sigma.n$ sudah digunakan}
    \\[1ex]
    \hline
    \iftag{prvBox}{
      \AxiomC{\sFmla{\True}{\Box !A}[\sigma.n]}
      \RightLabel{4r$\Box$}
      \UnaryInfC{\sFmla{\True}{\Box!A}[\sigma]}
      \DisplayProof}{}
    \iftag{notprvBox,notprvDiamond}{}{&}
    \iftag{prvDiamond}{
      \AxiomC{\sFmla{\False}{\Diamond !A}[\sigma.n]}
      \RightLabel{4r$\Diamond$}
      \UnaryInfC{\sFmla{\False}{\Diamond!A}[\sigma]}
      \DisplayProof}{}
    \\[1ex]
    \hline
    \end{tabular}
  \end{center}
  \caption{Aturan modal tambahan.}
  \ollabel{tab:more-rules}
\end{table}

Penambahan aturan-aturan ini menghasilkan sistem yang sahih dan lengkap
bagi logika yang diberikan dalam \olref{tab:logics-rules}.

\begin{table}
  \begin{center}
    \begin{tabular}{lll}
      \hline
      Logika & $R$ bersifat \dots & Aturan\\
      \hline
      $\Log{T} = \Log{KT}$ & refleksif &
      \iftag{prvBox}{T$\Box$}{}%
      \iftag{notprvBox,notprvDiamond}{}{, }%
      \iftag{prvDiamond}{T$\Diamond$}{}
      \\ \hline
      $\Log{D} = \Log{KD}$ & serial &
      \iftag{prvBox}{D$\Box$}{}%
      \iftag{notprvBox,notprvDiamond}{}{, }%
      \iftag{prvDiamond}{D$\Diamond$}{}
      \\ \hline
      $\Log{K4}$ & transitif &
      \iftag{prvBox}{4$\Box$}{}%
      \iftag{notprvBox,notprvDiamond}{}{, }%
      \iftag{prvDiamond}{4$\Diamond$}{}
      \\ \hline
      $\Log{B} = \Log{KTB}$ & refleksif, &
      \iftag{prvBox}{T$\Box$}{}%
      \iftag{notprvBox,notprvDiamond}{}{, }%
      \iftag{prvDiamond}{T$\Diamond$}{}\\
      & simetris &
      \iftag{prvBox}{B$\Box$}{}%
      \iftag{notprvBox,notprvDiamond}{}{, }%
      \iftag{prvDiamond}{B$\Diamond$}{}
      \\ \hline
      $\Log{S4} = \Log{KT4}$ & refleksif, &
      \iftag{prvBox}{T$\Box$}{}%
      \iftag{notprvBox,notprvDiamond}{}{, }%
      \iftag{prvDiamond}{T$\Diamond$}{},\\
      & transitif &
      \iftag{prvBox}{4$\Box$}{}%
      \iftag{notprvBox,notprvDiamond}{}{, }%
      \iftag{prvDiamond}{4$\Diamond$}{}
      \\ \hline
      $\Log{S5} = \Log{KT4B}$ & refleksif, &
      \iftag{prvBox}{T$\Box$}{}%
      \iftag{notprvBox,notprvDiamond}{}{, }%
      \iftag{prvDiamond}{T$\Diamond$}{},\\
      & transitif, &
      \iftag{prvBox}{4$\Box$}{}%
      \iftag{notprvBox,notprvDiamond}{}{, }%
      \iftag{prvDiamond}{4$\Diamond$}{},\\
      & Euklidean &
      \iftag{prvBox}{4r$\Box$}{}%
      \iftag{notprvBox,notprvDiamond}{}{, }%
      \iftag{prvDiamond}{4r$\Diamond$}{}
      \\ \hline
    \end{tabular}
  \end{center}
  \caption{Aturan !!^{tableau} bagi berbagai logika modal.}
  \ollabel{tab:logics-rules}
\end{table}


\begin{ex}
  Kita memberikan tableau tertutup yang menunjukkan bahwa
  $\Log{S5} \Proves \Ax{5}$, yakni $\Box!A \lif \Box\Diamond!A$.
  \begin{oltableau}
    [\pFmla{\False}{\Box\formula{A} \lif \Box\Diamond \formula{A}}{1},
      just = \TAss
      [\pFmla{\True}{\Box \formula{A}}{1}, just = {\TRule{\False}{\lif}[1]}
        [\pFmla{\False}{\Box\Diamond \formula{A}}{1},
          just = {\TRule{\False}{\lif}[1]}
          [\pFmla{\False}{\Diamond \formula{A}}{1.1},
            just = {\TRule{\False}{\Box}[3]}
            [\pFmla{\False}{\Diamond \formula{A}}{1},
              just = {4r$\Diamond$ 4}
              [\pFmla{\False}{\formula{A}}{1.1},
                just = {\TRule{\False}{\Diamond}[5]}
                [\pFmla{\True}{\formula{A}}{1.1},
                  just = {\TRule{\True}{\Box}[2]}, close]
              ]
            ]
          ]
        ]
      ]
    ]
  \end{oltableau}
\end{ex}

\begin{prob}
Berikan !!{tableau} tertutup yang menunjukkan hal-hal berikut:
  \begin{enumerate}
  \item $\Log{KT5} \Proves \Ax{B}$;
  \item $\Log{KT5} \Proves \Ax{4}$;
  \item $\Log{KDB4} \Proves \Ax{T}$;
  \item $\Log{KB4} \Proves \Ax{5}$;
  \item $\Log{KB5} \Proves \Ax{4}$;
  \item $\Log{KT} \Proves \Ax{D}$.
  \end{enumerate}
\end{prob}

\end{document}
